\section{Méthode}
\label{sec:methodology}

\subsection{Architectures}

Trois architectures sont comparées, toutes adaptées à la classification
sur les 200 classes de CUB-200-2011 :

\begin{itemize}
  \item \textbf{ResNet-50 pré-entraîné} : backbone \texttt{torchvision}
    pré-entraîné sur ImageNet, dont la dernière couche entièrement
    connectée est remplacée par une couche linéaire à 200 sorties.
  \item \textbf{ViT pré-entraîné} : backbone \texttt{timm}
    (\texttt{vit\_small\_patch16\_224} ou \texttt{vit\_small\_patch32\_224}
    selon la taille de patch), pré-entraîné sur ImageNet, avec une tête
    de classification à 200 sorties.
  \item \textbf{ViT from scratch} : implémentation entièrement codée pour
    ce projet, sans pré-entraînement. L'image est découpée en patches
    (16 ou 32 pixels) projetés linéairement (\emph{patch embedding}), un
    token \texttt{[CLS]} est ajouté à la séquence, un encodage
    positionnel appris est additionné, puis la séquence traverse 6
    blocs Transformer (dimension d'embedding 384, 6 têtes d'attention,
    ratio MLP de 4). La classification finale s'appuie sur le token
    \texttt{[CLS]} après normalisation.
\end{itemize}

\subsection{Protocole d'entraînement}

Chaque modèle est entraîné avec l'optimiseur AdamW (taux d'apprentissage
$10^{-4}$, weight decay $0{,}01$), une fonction de perte d'entropie
croisée, une taille de batch de 32, sur 3 epochs. Les images sont
prétraitées et augmentées selon les pipelines décrits en
Section~\ref{sec:data} (redimensionnement à 224$\times$224,
normalisation ImageNet, augmentation faible par défaut). Les
variantes ResNet-50 et ViT from scratch sont également entraînées sur
les sous-ensembles à 10\,\%, 25\,\% et 50\,\% du jeu d'entraînement
complet, dans les mêmes conditions, pour l'analyse de sensibilité à la
quantité de données (Section~\ref{sec:ablation}).

Le meilleur modèle de chaque run est sélectionné selon l'accuracy de
validation et sauvegardé automatiquement à chaque amélioration.

\subsection{Limite du protocole actuel}

Le nombre d'epochs (3) a été fixé pour des raisons de temps de calcul
disponible sur la durée du projet. Cette contrainte affecte
particulièrement le ViT from scratch, qui ne bénéficie d'aucun
pré-entraînement et nécessite un nombre d'itérations nettement plus
élevé pour converger, en cohérence avec son comportement data-hungry
documenté en Section~\ref{sec:related_work}. Un scheduler de taux
d'apprentissage, l'entraînement en précision mixte et un mécanisme
d'arrêt anticipé ont été implémentés dans l'outil d'entraînement du
projet mais n'ont pas encore été appliqués aux résultats rapportés
ici ; cette limite est reprise en Section~\ref{sec:discussion}.